\fancysection{Technical Approach and Methodology}
%
\label{sec:technical_approach}

To a large extent, our analysis plan is based on similar analysis from the
literature applied to \kepler{} and \tess{} lightcurves
\citep{Reinhold_et_al_13, Lurie_et_al_17, Windemuth_et_al_19,
CantoMartins_et_al_20, Justesen_Albrecht_21}, combining the best features of
each of these analysis and several other improvements of our own designed to
increasing the reliability of the catalog and allowing us to properly handle
difficulties with some EBs reported by these prior efforts. Our plans include
extensive validation procedures to ensure the resulting catalogs are as clean of
false positives or erroneously determined parameters as possible. Finally,
parameter uncertainties are a crucial piece of information, of almost equal
importance to the values of the parameters themselves. Our approach includes an
effort to thoroughly determine the uncertainties. Since observed properties
depend on complicated combinations of physical parameters, the uncertainty
distributions are highly non-Gaussian, and involve complicated correlations. We
will rely on a Markov chain Monte Carlo analysis to fully characterize these
uncertainty distributions, and the actual samples as well as summary statistics
for each EB in our catalog will be provided upon publication.

\subsection{Determining Eclipsing Binary Physical Parameters}
%
\label{sec:eb_parameters_technical}

\subsubsection{Lightcurve Preparation}
%
\label{sec:params_lc_preparation}

For a selected subset of stars \tess{} downloads pixel stamps with a 2-min
cadence and full images with a 30-min cadence. The stamps are processed by the
Science Processing Operations Center (SPOC) pipeline \citep{Jenkins_et_al_16}
providing  two classes of lightcurves: simple aperture photometry (SAP) and
pre-search data conditioning (PDC). The PDC lightcurves involve detrending
against empirically constructed co-trending basis vectors, which is specifically
designed to preserve astrophysical signals while removing systematics. While the
PDC corrections are not perfect, \citep{Justesen_Albrecht_21} found they are
well suited for high precision analysis of EBs. We will use the PDC lightcurves
by default, but for each target we also will begin by finding a maximum
likelihood model using our default modeling of PDC lightcurves and  a maximum
likelihood model modeling only the eclipses without the phase curve effects (see
below) using only the portions of the SAP lightcurves near eclipses (up to 2
times the eclipse duration) together with a second order polynomial model for
the out-of-eclipse variability. If the parameters of the two fits differ
substantially, the target will be flagged for manual inspection and potentially
re-analysis based on appropriately prepared SAP lightcurve (e.g. using the
\texttt{Lightkurve} package \citep{Lightkurve_18}).

For FFIs, we will use the Quick Look Pipeline (QLP) lightcurves
\citep{Huang_et_al_20a, Huang_et_al_20b, Kunimoto_et_al_21, Kunimoto_et_al_22}.
Those also come in two flavors: with and without detrending. The QLP authors
themselves caution that the detrending procedure will significantly affect low
frequency signals. Hence, we begin with non-detrended lightcurves, and adapt the
QLP detrending procedure for our application by modifying their high-pass filter
to preserve signals at or below the orbital frequency and applying the
detrending as part of MCMC after subtracting the EB model lightcurve (see
\refsec{params_bayesian_analysis}). To preserve modulations on the orbital
frequency or faster, we need to modify the minimum spacing allowed between the
smoothing spline nodes used for detrending to match the orbit. The rest of the
detrending procedure of the QLP will be preserved.

The first step in our analysis will be to select candidate EB lightcurves. For
stamp targets, \citep{Prsa_et_al_22} already conducted a thorough search,
combining multiple detection  methods, including two manual searches by citizen
scientists \citep{Eisner_et_al_21, Kristiansen_et_al_22}, a list of EB
candidates selected as \tess{} targets, and several eclipse or periodic signal
search pipelines \citep{Jenkins_et_al_16, Wheeler_Kipping_19, Kruse_et_al_19,
IJspeert_et_al_21} to identify 4584 EBs. These will constitute the list of
sources we will analyze from stamp data.

There is no equivalent catalog of EBs for FFI lightcurves. Instead, we will
follow \citep{Justesen_Albrecht_21} who searched the first year of \tess{} stamp
data for high signal-to-noise EBs. Their search was conducted in several steps.
First, a Lomb-Scargle periodogram was generated for each lightcurve, and stars
with no significant power were eliminated. Selected sources were subjected to
the transit least squares (TLS) algorithm developed by \citep{Hippke_Heller_19},
again excluding targets with no significant signal. Candidates for which neither
the eclipse depths, nor the depths of every other eclipse match within
5-$\sigma$ are also discarded as false positives. For the survivors, the
detected eclipses are masked and another search is conducted to check for
secondary eclipses of eccentric binaries which are also checked for consistency
of their depths. Finally, only high signal-to-noise candidates are kept (primary
eclipse > 5\% and secondary eclipse > 1\%). \citep{Justesen_Albrecht_21} report
that this procedure eliminated 98\% of the input sources. However, since we will
begin with $10^7$ lightcurves, $\sim2\times10^5$ will pass all tests. These are
too many to examine by eye, as \citep{Justesen_Albrecht_21} did. Instead, we
will find a maximum likelihood estimate for the system parameters and
automatically reject binaries for which the residuals between the observed
lightcurve and the model for the maximum likelihood parameters are too large.
The exact threshold for rejection will be determined by visual inspection of
several thousand examples. Since the likelihood function defined below is
non-linear and likely contains multiple local minima, we will use one of the
built-in global optimization algorithms of the scipy Python package. We will
experiment with basin hopping \citep{Wales_Doye_97}, differential evolution
\citep{Storn_Price_97}, and dual annealing \citep{Xiang_et_al_97} algorithms as
part of this project, and if no clear winner can be found we will use multiple
selecting whichever one produces the highest likelihood for each candidate. Even
after this step, we will end up with several tens of thousands of candidates.
All surviving EBs will be subjected to full analysis and included in the final
catalog, but we will manually vet the 5\,000 -- 10\,000 with tightest
constraints to produce a sub-catalog of highly reliable EBs.

\hl{For each lightcurve selected for full analysis, we will suppress variability
    on timescales longer than the orbital period (determined from the maximum
    likelihood optimization above). We will experiment with two approaches:
    median filter with window equal to the orbital period after masking the
    eclipses and folding the lightcurve at the orbital period and median
    filtering with a window width 1\% or less of the primary or secondary
    eclipse duration, whichever is smaller. The filtering and maximum likelihood
    will be applied iteratively to find an orbital period that maximizes the
likelihood of the filtered lightcurve.}

\subsubsection{Lightcurve and Spectral Energy Distribution Modeling Modeling}
%
\label{sec:params_lc_sed_modeling}

The most important constraint on the physical parameters of EBs comes from
comparing predicted lightcurves to what \tess{} observes. We require a model for
two types of effects: each star blocking some part of the light from its
companion from reaching us (i.e. the primary and secondary eclipses) and
out-of-eclipse variability, driven by interactions between the two stars (phase
curve). For the eclipses, we will use the Python package \texttt{batman}
\citep{Kreidberg_15}. We generate two \texttt{batman} models for a given set of
EB parameters, modeling the eclipse of the primary and secondary stars
separately. The two models are then added per the corresponding brightness ratio
of the stars.

Our model for the phase curve must handle ellipsoidal modulations, Doppler
beaming, and reflection effects. Let $\theta$ be the angle between the line of
sight direction and the line connecting the centers of the two stars in the
binary, and $d$ be the separation between the two stars. To an excellent
approximation, the ellipsoidal modulations are proportional to
$\sin(2\theta)/d^3$, the reflection effect is proportional to $\cos\theta/d^2$,
and the Doppler beaming is proportional to the line of sight velocity of the
primary star. The proportionality constants of these terms depend on details of
the stellar atmosphere and tidal love number. In our lightcurve model, we treat
these constants as model parameters free to be adjusted to match the lightcurve.
In fact in order to accommodate imperfectly filtered instrumental effects or
astrophysical variability not related to binary interactions, we also include
terms proportional to $\sin\theta/d^2$ and $\cos(2\theta)/d^3$.

It should be noted that while there are much more sophisticated treatments of
these effects available, for example by the \texttt{PHOEBE} package, those
necessarily are much more computationally intensive. For example, generating a
single EB lightcurve using \texttt{PHOEBE} with sufficient resolution for this
project  takes several to several tens of seconds using a single thread on a
contemporary computer. Given that \citep{Windemuth_et_al_19} found it takes
$\sim10^8$ samples ($10^6$ steps of 128 walker MCMC) for the samples to fully
converge to the desired posterior distribution, the proposed analysis is
completely infeasible for tens of thousands of EBs even using state of the art
high performance computing facilities. In contrast, the proposed modeling is
$10^3$ to $10^5$ times faster (depending on the binary parameters) making it
entirely feasible.

The apparent brightness of a system often places better constraints on the
masses than the lightcurve alone, because mass is tightly correlated with
luminosity. We rely on parallax based distances from Gaia Data Release 2
\citep{Bailer-Jones_et_al_18}. While Gaia DR3 parallaxes are already available
\citep{GaiaDR3_22}, converting to distance and uncertainty estimates is highly
non-trivial. If DR3 distance analysis becomes available before or during this
project, we will of course us the updated distances. However, given that we are
only analyzing very bright nearby stars, the DR2 precision is more than
sufficient for our purposes, and is not the dominant source of uncertainty when
converting absolute to apparent magnitudes by a significant margin. We will use
the extinction estimate fro the \tess{} input catalog (TIC) whenever available,
using the 3D galactic dust map by \citep{Green_et_al_19} as a backup, and the
dust map by \citep{Schlegel_et_al_98} as a last resort. For lightcurve modeling
we calculate difference it \tess{} magnitudes of the sources, and SED
constraints will be based on whichever of the following magnitudes are
available: Johnson–Cousins B and V, 2MASS J, H, K, nd WISE W1 and W2.

In order to build a model, we need the radii of the primary and secondary stars,
their extinction and reddening corrected apparent magnitudes, and a limb
darkening model. Without radial velocity observations, the masses of the stars
are not constrained, because the lightcurve depends only on orbital period,
eccentricity, and orientation, the ratio of radii, and the semimajor axis to
primary radius ratio. In order to compensate for the lack of radial velocity
data, we will instead rely on stellar evolution models to convert the two
stellar masses, and common age and metallicity for a binary to radii and
absolute magnitudes in the \tess{} bandpass for lightcurve modeling, and other
bandpasses for SED modeling (see above). In particular, we will use isochrones
generated using the Padova and Trieste Stellar Evolutionary Code
\citep{Bressan_et_al_12, Chen_et_al_14, Chen_et_al_15, Tang_et_al_14,
Marigo_et_al_17, Pastorelli_et_al_19, Pastorelli_et_al_20}, together with
bolometric corrections and star-by-star extinction coefficients (properly
accounting for the spectral dependence of extinction) as in
\citep{Chen_et_al_19}. \citep{CMD_interface} provides on-demand isochrones for a
range of stellar properties and we have already implemented high-precision
interpolation. We will assuming quadratic limb darkening, and treat the limb
darkening coefficients of each stars as model parameters to be constrained by
observations rather than using theoretical models.

\subsubsection{Bayesian Analysis}
%
\label{sec:params_bayesian_analysis}

\begin{wraptable}{l}{0.6\textwidth}
%
    \vspace{-\baselineskip}
%
    \caption{Model parameters sampled during MCMC}
%
    \label{tab:mcmc_params}
%
    \begin{tabular}{cl}
%
        \textbf{Parameter} & \textbf{Description}\\
%
        \toprule
%
        $M_1$, $M_2$ & Masses of the two stars\\
%
        $\tau, \left[\frac{Fe}{H}\right]$ & $\log_{10}$(Age) and metallicity of
        the binary\\
%
        $u_{ij},\ i,j\in\{1,2\}$ & Limb darkening coefficients\\
%
        $c_i, s_j,\ i,j\in\{1,2\}$ & Coefficients of the phase curve model\\
%
        $f_0$ & Flux zero-point\\
%
        $P_{orb}, e, i, \omega$ & \makecell{Orbital period, eccentricity,
        inclination,\\ and longitude of periapsis}\\
%
        d, E(B-V) & Distance to EB and reddening\\
%
        \makecell{$\sigma_{SED,sys}, \sigma_{LC,sys}$,\\$\sigma_d,
        \sigma_{E(B-V)}$} & Floating uncertainties (see text)
%
    \end{tabular}
%
\end{wraptable}
%
As described above, a full model of the lightcurves and SED requires specifying
23 parameters: Table \ref{tab:mcmc_params}.

\subsection{Physical Parameter Validation}
%
\label{sec:params_validation}

\subsection{Determining Spins}
%
\label{sec:spin_technical}

\subsubsection{Lightcurve Preparation}
%
\label{sec:spin_lc_preparation}

\subsection{Candidate Selection and Spin Period Measurements}
%
\label{sec:spin_search_and_measurement}

\subsection{Validating Measured Spins}
%
\label{sec:spin_validation}
